\documentclass{beamer}
\usetheme[faculty=ped,basePath=../fibeamer,logo=../fibeamer/logo/mu/Logo-UC-3.png]{fibeamer}
\graphicspath{{../resources/}}

\usepackage{algorithm}
\usepackage{algpseudocode}

\usetheme[faculty=ped]{fibeamer}
\usepackage[utf8]{inputenc}
\usepackage[
  main=english,
  czech, slovak
]{babel}
\title{Programación Orientada a Objetos}
\subtitle{Semestre II-2026}
\author{Profesora: Andreina Cota Vidaurre}

\usepackage{ragged2e}
\usepackage{booktabs}
\usepackage{tabularx}
\usepackage{tikz}
\usetikzlibrary{calc, shapes, backgrounds}
\usetikzlibrary{positioning, arrows.meta, shadows.blur, shapes.geometric}
\usepackage{amsmath, amssymb}
\usepackage{url}
\usepackage{listings}
\usepackage{graphicx}
\usepackage{amsmath}
\usepackage[T1]{fontenc}
\usepackage{xcolor}
\usepackage{minted}

\lstdefinelanguage{PseudocodigoES}{
    morekeywords={INICIO,FIN,PARA,HACER,SI,ENTONCES,FIN_SI,FIN_PARA,IMPRIMIR},
    sensitive=true,
    morecomment=[l]{//},
    morestring=[b]"
}

\lstdefinelanguage{PseudocodeEN}{
    morekeywords={BEGIN,END,FOR,DO,IF,THEN,END_IF,END_FOR,PRINT},
    sensitive=true,
    morecomment=[l]{//},
    morestring=[b]"
}

\lstset{
    basicstyle=\ttfamily\small,
    keywordstyle=\bfseries,
    columns=fullflexible,
    keepspaces=true,
    showstringspaces=false,
    frame=single,
    xleftmargin=0em,
    framexleftmargin=0em,
    framesep=2pt,
    gobble=4,
    resetmargins=true,
    aboveskip=0pt,
    belowskip=0pt
}

\lstdefinestyle{pythonstyle}{
    language=Python,
    basicstyle=\ttfamily\small,
    keywordstyle=\bfseries,
    commentstyle=\itshape,
    stringstyle=\ttfamily,
    showstringspaces=false,
    columns=fullflexible,
    keepspaces=true,
    frame=single,
    breaklines=true,
    xleftmargin=0em,
    framexleftmargin=0em,
    framesep=2pt,
    aboveskip=0pt,
    belowskip=0pt,
    gobble=4,
    resetmargins=true,
    emph={print,input,int,float,str,bool,len,range},
    emphstyle=\bfseries
}

\lstset{style=pythonstyle}

\tikzstyle{iniciofin} = [
    ellipse,
    minimum width=2.5cm,
    minimum height=1cm,
    text centered,
    draw=black,
    fill=blue!25
]

\tikzstyle{proceso} = [
    rectangle,
    rounded corners,
    minimum width=3cm,
    minimum height=1cm,
    text centered,
    draw=black,
    fill=cyan!20
]

\tikzstyle{decision} = [
    diamond,
    aspect=2,
    text centered,
    draw=black,
    fill=yellow!35
]

\tikzstyle{flecha} = [
    thick,
    ->,
    >=Stealth
]

\frenchspacing
\begin{document}
  \shorthandoff{-}
  \frame[c]{\maketitle}

    %------------------------------------------------------------------
    \begin{frame}{Ejercicio: Autonomía de un bus eléctrico}
        En los sistemas de transporte público con electromovilidad, la planificación
        de recorridos requiere estimar si la energía almacenada en la batería de un
        bus es suficiente para completar una ruta.

        \vspace{0.4em}
        Cada tramo de la ruta tiene un consumo energético conocido, determinado por
        su distancia, pendiente y condiciones de tránsito.

        \vspace{0.6em}
        \textcolor{blue}{\textbf{Tu tarea:} modelar la batería de un bus eléctrico
        mediante una clase y utilizarla para simular el recorrido de una ruta.}
    \end{frame}

    \begin{frame}[fragile, t]{Parte 1: La clase \texttt{Bateria}}
        Toda batería se identifica por un \textbf{código}, tiene una
        \textbf{capacidad máxima} y una \textbf{carga actual}, ambas medidas en kWh.

        \vspace{0.4em}
        Define la clase \mintinline{python}{Bateria} construyendo los siguientes métodos:

        \begin{enumerate}
            \item \textbf{Método de inicialización}
            \mintinline{python}{__init__(self, id_bateria, capacidad, carga)}:
            crea la batería y define sus tres atributos a partir de los valores recibidos.
            \pause
            \item \textbf{Método} \mintinline{python}{descargar(self, energia)}:
            extrae de la batería la energía solicitada. Si no tiene suficiente carga,
            entrega solo la disponible. La carga actual nunca puede quedar negativa.
            Debe \textbf{retornar la energía que efectivamente entregó}.
            \pause
            \item \textbf{Método} \mintinline{python}{estado(self)}:
            retorna un texto con el formato \texttt{"BUS-E12: 120/350 kWh"},
            es decir, código, carga actual y capacidad.
        \end{enumerate}
    \end{frame}

    \begin{frame}[fragile, t]{Parte 2: Simulación de la ruta}
        Un bus inicia su recorrido con una batería de \textbf{350 kWh} de capacidad
        y \textbf{120 kWh} de carga. La ruta está representada como una lista de tuplas
        \texttt{(nombre del tramo, consumo en kWh)}:
        \begin{minted}[xleftmargin=-2cm, frame=leftline, fontsize=\small]{python}
            tramos = [
                ("Terminal - Centro", 30),
                ("Centro - Hospital", 45),
                ("Hospital - Universidad", 25),
                ("Universidad - Terminal", 40)
            ]
        \end{minted}
    \end{frame}

    \begin{frame}[fragile, t]{Parte 2: Simulación de la ruta}
        Escribe un programa que:
        \begin{enumerate}
            \item Cree un objeto de la clase \mintinline{python}{Bateria} con los datos
            del bus y muestre su estado inicial.
            \pause
            \item Recorra los tramos en orden y, para cada uno, descargue de la batería
            el consumo correspondiente.
            \pause
            \item Muestre un mensaje por cada tramo:
            \begin{itemize}
                \item Si entregó toda la energía:
                \texttt{"<tramo>: completado | <estado>"}
                \item Si no alcanzó:
                \texttt{"<tramo>: energía insuficiente, faltaron <faltante> kWh"}
            \end{itemize}
        \end{enumerate}

        \vspace{0.5em}
    \end{frame}

    \begin{frame}[fragile, t]{Plantilla: clase \texttt{Bateria}}
        Complete los métodos de la clase:
        \begin{minted}[xleftmargin=-2cm, frame=leftline, fontsize=\scriptsize]{python}
            class Bateria:
                def __init__(self, id_bateria, capacidad, carga):
                    # Completar
                    pass

                def descargar(self, energia):
                    # Completar
                    pass

                def estado(self):
                    # Completar
                    pass
        \end{minted}
    \end{frame}

    \begin{frame}[fragile, t]{Plantilla: programa principal}
        Complete la creación del objeto y el recorrido de los tramos:
        \begin{minted}[xleftmargin=-2cm, frame=leftline, fontsize=\scriptsize]{python}
            # Programa principal
            tramos = [
                ("Terminal - Centro", 30),
                ("Centro - Hospital", 45),
                ("Hospital - Universidad", 25),
                ("Universidad - Terminal", 40)
            ]

            bus = # Completar: crear la batería del bus
            print(bus.estado())

            for tramo, consumo in tramos:
                # Completar
                pass
        \end{minted}
    \end{frame}

    \begin{frame}[fragile, t]{Salida esperada}
        \begin{minted}[xleftmargin=-2cm, frame=leftline, fontsize=\small]{text}
            BUS-E12: 120/350 kWh
            Terminal - Centro: completado | BUS-E12: 90/350 kWh
            Centro - Hospital: completado | BUS-E12: 45/350 kWh
            Hospital - Universidad: completado | BUS-E12: 20/350 kWh
            Universidad - Terminal: energía insuficiente, faltaron 20 kWh
        \end{minted}
    \end{frame}

\end{document}
